\section{Bipartisanship: can one caucus support it without losing the other?}
\label{sec:bipartisanship}

The sixth question decides whether a bill is a message or a law. A purely partisan
bill dies in the other chamber; H. R. 9510 is unusual because the same provision
answers two caucuses that start from different premises.

To a member focused on innovation and a lighter regulatory hand, verification before
generation is a predictable, automatable rule that lets credible systems deploy
faster, replaces slow and inconsistent manual review with a fixed gate, and creates
no new entitlement, tax, or fee. That posture is consistent with the executive
direction to remove barriers to American leadership in artificial intelligence
\cite{eo14179}. To a member focused on patient safety and human oversight, the same
section is a hard safety floor: a ten-gate examination with three catastrophe
predicates that cannot be waived, a preserved human-over-AI role, and a tamper-evident
record an auditor can read. Neither side has to adopt the other's reasoning to reach
the same vote. \autoref{fig:converge} shows the convergence.

\begin{psgfig}
\node[psgevid] (i) at (0,1.1) {Innovation and\\deregulation appeal};
\node[psgact] (s) at (0,-1.1) {Patient safety and\\oversight appeal};
\node[psgtwo] (m1) at (4.2,1.1) {Faster credible\\deployment};
\node[psgtwo] (m2) at (4.2,-1.1) {Hard floor,\\no entitlement};
\node[psgact] (y) at (8.4,0) {One yes vote};
\draw[psgedge] (i)--(m1); \draw[psgedge] (s)--(m2);
\draw[psgedge] (m1) to[out=0,in=120,looseness=0.7] (y.north west);
\draw[psgedge] (m2) to[out=0,in=-120,looseness=0.7] (y.south west);
\end{psgfig}
\psgcaption{\textbf{Figure 7.} Two appeals, one vote: an innovation rationale and a
safety rationale converge on the same provision; the colored TikZ reproduction of
catalog diagram 12. Both curved arrows carry an explicit looseness of 0.7.}

The political science of getting a bill across the aisle supports the approach.
Legislators move when a case is framed for their own priorities and paired with
credible evidence rather than advocacy alone, and the same evidence can be read by
both sides when ambiguity, not just uncertainty, is reduced
\cite{moreland2015hearing, brownson2009researchers, cairney2016evidence}. The
framework's discipline of pairing every claim with a cited fact is built for exactly
that.

Bipartisanship is answered by design: a single mechanism that a deregulation-minded
member and a safety-minded member can each support on their own terms, with no
daylight that an opponent can exploit.
